% PUC-SP e PUC-Camp 1995 — 19 questões MCQ — História

---
tipo: Múltipla Escolha
dificuldade: Média
nivel: médio
disciplina: História
assunto: Grandes Navegações
gabarito: A
tags: [expansão marítima, Portugal, Fernando Pessoa, Era Moderna]
fonte: concurso
concurso: PUC
banca: PUC
ano: 1995
numero: 1
---
\question
``Ó mar salgado, quanto do teu sal / São lágrimas de Portugal! / Por te cruzarmos, quantas mães choraram, / Quantos filhos em vão rezaram! / Quantas noivas ficaram sem casar / Para que fosses nosso, ó mar! / Valeu a pena? Tudo vale a pena / Se a alma não é pequena. / Quem quer passar além do Bojador / Tem que passar além da dor. / Deus ao mar o perigo e o abismo deu. / Mas nele é que espelhou o céu.'' (Fernando Pessoa, \textit{Mar Português}.)

O poema refere-se à conquista dos mares pelos portugueses, o início da era moderna. Se os resultados mais conhecidos das Grandes Navegações foram a abertura de novas rotas comerciais, a conquista de novas terras e o espalhamento da cultura europeia, alguns dos elementos cujo contexto histórico auxilia na compreensão das origens dessa expansão marítima são:
\begin{choices}
\CorrectChoice o avanço das técnicas de navegação; a busca do mítico paraíso terrestre; a percepção do universo segundo uma ordem racional.
\choice o mito do abismo do mar; a desmonetarização da economia; a vontade de enriquecimento rápido.
\choice a busca de ouro para as Cruzadas; a descentralização monárquica; o desenvolvimento da matemática.
\choice a demanda de especiarias; a aliança com as cidades italianas; a ânsia de expandir o cristianismo.
\choice o anseio de crescimento mercantil; os relatos de viajantes medievais; a conquista de Portugal pelos mouros.
\end{choices}

---
tipo: Múltipla Escolha
dificuldade: Média
nivel: médio
disciplina: História
assunto: Era Vargas e Estado Novo
gabarito: B
tags: [CLT, trabalhismo, sindicatos, Vargas, leis trabalhistas]
fonte: concurso
concurso: PUC
banca: PUC
ano: 1995
numero: 2
---
\question
No Brasil, a CLT --- Consolidação das Leis do Trabalho --- foi criada pelo Decreto 5452, de 1943, em meio ao governo de Getúlio Vargas, para reunir e sistematizar as leis trabalhistas existentes no país. Tais leis representaram a:
\begin{choices}
\choice conquista evidente do movimento operário sindical e partidariamente organizado desde 1917, defensor de projetos socialistas e responsável pela ascensão de Vargas ao poder.
\CorrectChoice participação do Estado como árbitro na mediação das relações entre patrões e trabalhadores de 1930 em diante, permitindo a Vargas propor a racionalização e a despolitização das reivindicações trabalhistas.
\choice inspiração notadamente fascista, que orientou o Estado Novo desde sua implantação em 1937, desviando Vargas das intenções nacionalistas presentes no início de seu governo.
\choice atuação controladora do Estado brasileiro sobre os sindicatos e associações de trabalhadores, permitindo a Vargas criar, a partir de 1934, o primeiro partido político de massas da história brasileira.
\choice pressão norte-americana, que se tornou mais clara após 1945, para que Vargas controlasse os grupos anárquicos e socialistas presentes nos movimentos operário e camponês.
\end{choices}

---
tipo: Múltipla Escolha
dificuldade: Difícil
nivel: médio
disciplina: História
assunto: Roma Antiga
gabarito: C
tags: [Cristianismo, Estado Romano, Igreja, século IV, aliança]
fonte: concurso
concurso: PUC
banca: PUC
ano: 1995
numero: 3
---
\question
Dentre os itens a seguir, dois representam características integrantes do ideário cristão que, à época do reconhecimento do Cristianismo como religião oficial de Roma (séc. IV), funcionaram como elementos facilitadores da aliança entre a Igreja Cristã e o Estado Romano:
\begin{enumerate}
\item o dogma da transcendência divina.
\item as noções de culpa original dos homens e de perdão divino.
\item os dogmas da criação e do juízo final.
\item o missionarismo expansionista.
\item a moral celibatária.
\item as concepções de inferno, purgatório e reino dos céus.
\item a estrutura hierárquica da organização clerical.
\end{enumerate}
Os itens corretos são os de número:
\begin{choices}
\choice 5 e 1.
\choice 3 e 6.
\CorrectChoice 4 e 7.
\choice 6 e 4.
\choice 3 e 7.
\end{choices}

---
tipo: Múltipla Escolha
dificuldade: Difícil
nivel: médio
disciplina: História
assunto: Revolução Francesa
gabarito: E
tags: [revolução burguesa, burguesia, antigo regime, república, França]
fonte: concurso
concurso: PUC
banca: PUC
ano: 1995
numero: 4
---
\question
O motivo pelo qual o conjunto de mudanças políticas que resultou na implantação do regime republicano na França, no século XVIII, pode ser classificado como uma revolução burguesa, é o fato de que nesse processo:
\begin{choices}
\choice a estrutura social francesa viu-se reduzida a uma polarização entre o bloco de apoio ao antigo regime --- aristocracia, camponeses e trabalhadores urbanos --- e o bloco de apoio à república operário-burguesa.
\choice a burguesia conseguiu a adesão ideológica da aristocracia, especialmente quanto à ``abertura das carreiras públicas aos talentos individuais'', possibilitando a ascensão de seus representantes ao poder do Estado.
\choice o comando da burguesia desde o início se revelou irrefutável, colocando a serviço de seus objetivos revolucionários os mais variados setores da população, liderando assim uma restauração do Antigo Regime.
\choice as vanguardas operário-camponesas se colocaram ao lado da burguesia, pois tinham claro que suas reivindicações somente alcançariam consequência numa sociedade com relações burguesas de produção já desenvolvidas.
\CorrectChoice os resultados políticos das sucessivas convulsões sociais geradas nos quadros da crise do estado monárquico francês foram, ao final, capitalizados pela burguesia, que pôde assim dar início à viabilização de seus interesses políticos e econômicos.
\end{choices}

---
tipo: Múltipla Escolha
dificuldade: Média
nivel: médio
disciplina: História
assunto: Era Napoleônica
gabarito: A
tags: [Napoleão, expansão, liberalismo, nacionalismo, Europa]
fonte: concurso
concurso: PUC
banca: PUC
ano: 1995
numero: 5
---
\question
Relativamente à expansão napoleônica (1805--1815), pode-se afirmar que acarretou mudanças no quadro político europeu, tais como:
\begin{choices}
\CorrectChoice difusão do ideal revolucionário liberal, ampliação temporária do raio de influência francesa e fortalecimento do ideário nacionalista nos países dominados.
\choice isolamento diplomático da nação inglesa, radicação definitiva do republicanismo no continente e estabelecimento do equilíbrio geopolítico entre os países atingidos.
\choice desestabilização das monarquias absolutistas, estímulo para o desenvolvimento industrial nas colônias espanholas e implantação do belicismo entre as nações.
\choice desenvolvimento do cosmopolitismo entre os povos do império francês, incrementação da economia nos países ibéricos e contenção das lutas sociais.
\choice difusão do militarismo como forma de controle político, abertura definitiva do mercado mundial para os franceses e estímulo decisivo para as lutas anti-colonialistas.
\end{choices}

---
tipo: Múltipla Escolha
dificuldade: Média
nivel: médio
disciplina: História
assunto: Nova República e Governo Collor
gabarito: D
tags: [impeachment, Collor, FHC, Lula, partidos políticos, eleições 1994]
fonte: concurso
concurso: PUC
banca: PUC
ano: 1995
numero: 6
---
\question
As eleições presidenciais brasileiras de 1994 envolveram oito candidatos concorrendo por partidos ou alianças diversas. Alguns fizeram referências a episódios ou personagens da história política brasileira do século XX. Entre tais referências pode-se mencionar a lembrança do:
\begin{choices}
\choice nascimento de vários partidos entre 1979 e 1982, momento da ``reforma partidária'', quando surgiram, entre outros, o PMDB de Orestes Quércia e o PRN de Carlos Gomes.
\choice golpe militar de 1964, defendido àquela época pelo PFL e pelo PSC, que instalou no poder o almirante Fortuna, presidenciável nas últimas eleições.
\choice desenvolvimentismo de Juscelino Kubitschek, que governou de 1956 a 1961, e que fez aparecer a proposta social-democrata, defendida por Esperidião Amin e por Enéas Carneiro.
\CorrectChoice ``impeachment'' do presidente Fernando Collor de Mello, ocorrido em 1992 e que contou com a participação favorável, entre outros, de Luís Inácio Lula da Silva e de Fernando Henrique Cardoso.
\choice período presidencial de Getúlio Vargas entre 1951 e 1954, quando se formaram os atuais partidos políticos de esquerda --- PT e PTB --- e quando surgiu a liderança política de Leonel Brizola.
\end{choices}

---
tipo: Múltipla Escolha
dificuldade: Média
nivel: médio
disciplina: História
assunto: Feudalismo
gabarito: E
tags: [cavaleiros, nobreza, feudalismo, guerra, Marc Bloch]
fonte: concurso
concurso: PUC
banca: PUC
ano: 1995
numero: 7
---
\question
``A própria vocação do nobre lhe proibia qualquer atividade econômica direta. Ele pertencia de corpo e alma à sua função própria: a do guerreiro. (...) um corpo ágil e musculoso não é o bastante para fazer o cavaleiro ideal. É preciso ainda acrescentar a coragem. E é também porque proporciona a esta virtude a ocasião de se manifestar que a guerra põe tanta alegria no coração dos homens, para os quais a audácia e o desprezo da morte são, de algum modo, valores profissionais.'' (Marc Bloch, \textit{A Sociedade Feudal}.)

O autor fala da condição social dos nobres medievais e dos valores ligados às suas ações guerreiras. É possível dizer que a atuação guerreira desses cavaleiros representa, respectivamente, para a sociedade e para eles próprios:
\begin{choices}
\choice a garantia de segurança, um contexto em que as classes e os estados nacionais se encontram em conflito, e a perspectiva de conquistas de terras e riquezas.
\choice o cumprimento das obrigações senhoriais ligadas à produção, e a proibição da transmissão hereditária das conquistas realizadas.
\choice a permissão real para realização de atividades comerciais, e a eliminação do tédio de um cotidiano de cultura rudimentar e alheio a assuntos administrativos.
\choice o respeito às relações de vassalagem travadas entre senhores e servos, e a diversão sob a forma de torneios e jogos em épocas de paz.
\CorrectChoice a participação nas guerras santas e na defesa do catolicismo, e a possibilidade de pilhagem de homens e coisas, de massacres e mutilações de inimigos.
\end{choices}

---
tipo: Múltipla Escolha
dificuldade: Média
nivel: médio
disciplina: História
assunto: Entradas e Bandeiras
gabarito: D
tags: [bandeirantes, Brasil colonial, São Paulo, economia colonial]
fonte: concurso
concurso: PUC
banca: PUC
ano: 1995
numero: 8
---
\question
Personagem atuante no Brasil colônia, foi ``fruto social de uma região marginalizada, de escassos recursos materiais e de vida econômica restrita (...)'', teve suas ações orientadas ``ou no sentido de tirar o máximo proveito das brechas que a economia colonial eventualmente oferecia para a efetivação de lucros rápidos e passageiros em conjunturas favoráveis --- como no caso da caça ao índio --- ou no sentido de buscar alternativas econômicas fora do quadro da agricultura voltada para o mercado externo (...)''. (Carlos Henrique Davidoff, 1982.)

O personagem e a região a que o texto se refere são, respectivamente:
\begin{choices}
\choice o jesuíta e a província Cisplatina.
\choice o tropeiro e o vale do Paraíba.
\choice o caipira e o interior paulista.
\CorrectChoice o bandeirante e a província de São Paulo.
\choice o caiçara e o litoral baiano.
\end{choices}

---
tipo: Múltipla Escolha
dificuldade: Média
nivel: médio
disciplina: História
assunto: Revolução Russa e Stalinismo
gabarito: C
tags: [realismo socialista, URSS, arte oficial, cultura soviética, propaganda]
fonte: concurso
concurso: PUC
banca: PUC
ano: 1995
numero: 9
---
\question
O Estado Soviético, formado após a Revolução Russa, cuidou de expurgar da cultura do país toda manifestação artística associada ao chamado ``espírito burguês''. Foi criada uma política cultural que decretava como arte oficial apenas as expressões que servissem de estímulo para a ideologia do proletariado. Dessa forma, foi consagrado um estilo conhecido por:
\begin{choices}
\choice expressionismo soviético --- que, através de uma orientação estética intimista, procurava expor a ``alma inquieta dos povos eslavos'' que passaram a integrar a URSS.
\choice abstracionismo proletário --- que, através da decomposição geométrica do real, exprimia a ``ordenação sincrônica da sociedade comunista''.
\CorrectChoice realismo socialista --- que, através de composições didáticas, esteticamente simplificadas, procurava enaltecer a ``combatividade, a capacidade de trabalho e a consciência social'' do povo soviético.
\choice romantismo comunista --- que, através de um figurativismo apenas sugestivo, procurava realizar a ``idealização do mujique'', o camponês russo típico, como representante das raízes culturais russas.
\choice concretismo operário --- que, através de uma concepção criadora autônoma, utilizava elementos visuais e táteis, com o objetivo de mostrar a ``prevalência do concreto sobre o abstrato''.
\end{choices}

---
tipo: Múltipla Escolha
dificuldade: Difícil
nivel: médio
disciplina: História
assunto: Segundo Reinado
gabarito: B
tags: [D. Pedro II, poder moderador, burocracia, II Império, J. Murilo de Carvalho]
fonte: concurso
concurso: PUC
banca: PUC
ano: 1995
numero: 10
---
\question
``A enorme visibilidade do poder era sem dúvida em parte devida à própria monarquia com suas pompas, seus rituais, com o carisma da figura real. Mas era também fruto da centralização política do Estado. (...) este poder era em boa parte ilusório. A burocracia do Estado era macrocefálica: tinha cabeça grande mas braços muito curtos. Agigantava-se na corte mas não alcançava as municipalidades e mal atingia as províncias. (...) Daí a observação de que, apesar de suas limitações no que se referia à formulação e implementação de políticas, o governo passava a imagem do todo-poderoso, era visto como o responsável por todo o bem e todo o mal do Império.'' (J. Murilo de Carvalho, \textit{Teatro de Sombras}.)

O fragmento refere-se ao II Império brasileiro, controlado por D. Pedro II (1840--1889). Do ponto de vista político, o II Império pode ser representado como:
\begin{choices}
\choice palco de enfrentamento entre liberais e conservadores que, partindo de princípios políticos e ideológicos opostos, questionaram com igual violência essa aparente centralização e se uniram no Golpe da Maioridade.
\CorrectChoice jogo de aparências, em que a atuação política do Imperador conheceu as mudanças e os momentos de indefinição referidos --- refletindo as oscilações dos setores sociais hegemônicos ---, como bem exemplificado na questão da Abolição.
\choice cenário de várias revoltas de caráter regionalista --- entre elas a Farroupilha e a Cabanagem --- devido à incapacidade do governo imperial de controlar as províncias mais distantes da capital.
\choice universo de plena difusão das ideias liberais, o que implicou uma aceitação por parte do Imperador da diminuição de seus poderes, oferecendo condições para a proclamação da República.
\choice teatro para a plena manifestação do poder moderador que, desde a Constituição de 1824, permitia amplas possibilidades de intervenção para o Imperador e foi usado no Segundo Reinado para encerrar os conflitos entre liberais e socialistas.
\end{choices}

---
tipo: Múltipla Escolha
dificuldade: Fácil
nivel: médio
disciplina: História
assunto: Roma Antiga
gabarito: C
tags: [etruscos, Itália, povos primitivos, italiotas, gregos]
fonte: concurso
concurso: PUC
banca: PUC
ano: 1995
numero: 11
---
\question
Sobre os primitivos habitantes da Itália, pode-se afirmar que os:
\begin{choices}
\choice italiotas acomodaram-se no Sul da Itália, onde desenvolveram povoados.
\choice gregos ocuparam a parte central da Península, subdividindo-se em vários clãs.
\CorrectChoice etruscos, provavelmente originários da Ásia, ocuparam o Norte da Península.
\choice lígures fixaram-se ao Sul combatendo ferrenhamente os etruscos.
\choice sículos penetraram na Península através da cadeia dos Alpes e ocuparam o Norte.
\end{choices}

---
tipo: Múltipla Escolha
dificuldade: Média
nivel: médio
disciplina: História
assunto: Igreja e Papado na Idade Média
gabarito: B
tags: [Igreja, feudalismo, mosteiros, cultura, monopólio cultural]
fonte: concurso
concurso: PUC
banca: PUC
ano: 1995
numero: 12
---
\question
A Igreja integrou-se ao Sistema Feudal através dos mosteiros, cujas características se assemelhavam às dos domínios dos senhores feudais. Como tinha:
\begin{choices}
\choice o controle do destino espiritual, procurou combater a usura entre os integrantes do clero e entre os judeus, no que foi rigorosamente obedecida.
\CorrectChoice o monopólio da cultura, tinha também o monopólio da interpretação da realidade social.
\choice grande influência na formação da mentalidade, insistia no ideal do preço justo, permitindo que na venda dos produtos se cobrasse a mais apenas o custo do transporte.
\choice o controle da realidade social, exigia que os cristãos distribuíssem os excedentes entre seus parentes mais próximos para auferir lucros.
\choice a fiscalização sobre a distribuição dos excedentes em épocas de calamidade, inibia a atuação dos comerciantes inescrupulosos, ameaçando-os com multas ou com a perda de suas propriedades.
\end{choices}

---
tipo: Múltipla Escolha
dificuldade: Média
nivel: médio
disciplina: História
assunto: Bizâncio e Mundo Árabe
gabarito: D
tags: [Maomé, islamismo, Arábia, unificação religiosa, pré-islâmico]
fonte: concurso
concurso: PUC
banca: PUC
ano: 1995
numero: 13
---
\question
Para compreender a unificação religiosa e política da Arábia por Maomé, é necessário conhecer:
\begin{choices}
\choice a atuação das seitas religiosas sunita e xiita, que contribuíram para a consolidação do Estado teocrático islâmico.
\choice os princípios legitimistas obedecidos pela tribo coraixita, da qual fazia parte.
\choice os fundamentos do sincretismo religioso que marcou a doutrina islâmica.
\CorrectChoice as particularidades da vida dos árabes nos séculos anteriores ao surgimento do islamismo.
\choice a atuação da dinastia dos Omíadas que, misturando-se com os habitantes da região do Maghreb, converteram-se à religião muçulmana e passaram a ser chamados de mouros.
\end{choices}

---
tipo: Múltipla Escolha
dificuldade: Média
nivel: médio
disciplina: História
assunto: Absolutismo Monárquico
gabarito: C
tags: [Estado Moderno, absolutismo, mercantilismo, monarquia]
fonte: concurso
concurso: PUC
banca: PUC
ano: 1995
numero: 14
---
\question
Dentre as instituições políticas do Estado Moderno, aquela que mais o caracteriza é o:
\begin{choices}
\choice absolutismo monárquico, nova forma política cujos fundamentos estavam expressos na \textit{Suma Teológica} de Tomás de Aquino.
\choice mercantilismo que servia para justificar o enriquecimento da Igreja Católica, mas não traduzia os interesses do monarca absolutista.
\CorrectChoice absolutismo monárquico que intervinha na vida econômica.
\choice liberalismo praticado pelos Príncipes, mas limitado pela tradição e pelo equilíbrio entre as classes sociais.
\choice absolutismo monárquico que punha em prática uma política econômica de características não intervencionistas, quase liberais --- a política mercantilista.
\end{choices}

---
tipo: Múltipla Escolha
dificuldade: Fácil
nivel: médio
disciplina: História
assunto: Ciclo do Pau-Brasil
gabarito: A
tags: [pau-brasil, escambo, indígenas, Brasil colonial, pré-colonial, comércio]
fonte: concurso
concurso: PUC
banca: PUC
ano: 1995
numero: 15
---
\question
Em razão de as comunidades primitivas indígenas representarem, no Período Colonial, apenas reservas de força de trabalho a ser aproveitada no corte e transporte do pau-brasil, entre 1500 e 1530, no Brasil:
\begin{choices}
\CorrectChoice o comércio realizava-se através da troca direta ou escambo.
\choice a maioria das atividades produtivas concentrava-se na economia informal.
\choice o extrativismo mineral acabou desenvolvendo um mercado de consumo interno.
\choice a economia baseou-se essencialmente em atividades agrícolas.
\choice a expansão da pecuária impulsionou a utilização da mão-de-obra escrava africana.
\end{choices}

---
tipo: Múltipla Escolha
dificuldade: Média
nivel: médio
disciplina: História
assunto: Reforma Protestante e Contrarreforma
gabarito: E
tags: [Calvino, Calvinismo, Reforma Protestante, protestantismo, capital comercial]
fonte: concurso
concurso: PUC
banca: PUC
ano: 1995
numero: 16
---
\question
O Calvinismo foi:
\begin{choices}
\choice a doutrina que sintetizou as ideias dos reformadores que a antecederam, formulando o campo protestante em torno dos princípios do cesaropapismo e do culto dos santos.
\choice apenas um prolongamento das ideias preconizadas por Lutero, que admitia que o Príncipe, além de exercer poder civil absoluto, devia vigiar e governar, por direito divino, a Igreja cristã.
\choice um movimento originário na Suíça, como resultado de convulsões sociais locais, que revelavam uma manifestação de rebeldia contra as taxas cobradas pela Igreja e sobre a liberação da prática do divórcio.
\choice o resultado das preocupações pessoais de Ulrico Zuínglio e dos problemas relacionados com o celibato clerical.
\CorrectChoice a mais extremada seita protestante em relação ao Catolicismo e a mais próxima das questões levantadas, em termos éticos, pelo rápido desenvolvimento do capital comercial e financeiro.
\end{choices}

---
tipo: Múltipla Escolha
dificuldade: Fácil
nivel: médio
disciplina: História
assunto: Revolução Industrial
gabarito: A
tags: [proletariado, trabalho assalariado, meios de produção, Revolução Industrial]
fonte: concurso
concurso: PUC
banca: PUC
ano: 1995
numero: 17
---
\question
Dentre as consequências sociais forjadas pela Revolução Industrial pode-se mencionar:
\begin{choices}
\CorrectChoice o desenvolvimento de uma camada social de trabalhadores que, destituídos dos meios de produção, passaram a sobreviver apenas da venda de sua força de trabalho.
\choice a melhoria das condições de habitação e sobrevivência para o operariado, proporcionada pelo surto de desenvolvimento econômico.
\choice a ascensão social dos artesãos que reuniram seus capitais e suas ferramentas em oficinas ou domicílios rurais dispersos, aumentando os núcleos domésticos de produção.
\choice a criação do Banco da Inglaterra, com o objetivo de financiar a monarquia e ser também uma instituição geradora de empregos.
\choice o desenvolvimento de indústrias petroquímicas favorecendo a organização do mercado de trabalho, de maneira a assegurar emprego a todos os assalariados.
\end{choices}

---
tipo: Múltipla Escolha
dificuldade: Fácil
nivel: médio
disciplina: História
assunto: Revolução Francesa
gabarito: B
tags: [Terror, Comitê de Salvação Pública, Convenção, Robespierre, Revolução Francesa]
fonte: concurso
concurso: PUC
banca: PUC
ano: 1995
numero: 18
---
\question
No contexto da Revolução Francesa, a organização do Governo Revolucionário significou uma forte centralização do poder: o Comitê de Salvação Pública, eleito pela Convenção, passou a ser o efetivo órgão do Governo. Havia ainda o Comitê de Segurança Geral, que dirigia a polícia e a justiça, subordinado ao Tribunal Revolucionário, que tinha competência para punir até a morte todos os suspeitos de oposição ao regime. O conjunto de medidas de exceção adotadas pelo Governo revolucionário deram margem a que essa fase da Revolução viesse a ser conhecida como:
\begin{choices}
\choice os Massacres de Setembro.
\CorrectChoice o Período do Terror.
\choice o Grande Medo.
\choice o Período do Termidor.
\choice o Golpe do 18 de Brumário.
\end{choices}

---
tipo: Múltipla Escolha
dificuldade: Média
nivel: médio
disciplina: História
assunto: Independência do Brasil
gabarito: C
tags: [abertura dos portos, colonialismo, influência britânica, independência, D. João VI]
fonte: concurso
concurso: PUC
banca: PUC
ano: 1995
numero: 19
---
\question
A abertura dos portos teve um alcance histórico profundo, pois deu início a um duplo processo: o:
\begin{choices}
\choice do desenvolvimento do primeiro surto manufatureiro no Brasil e o crescimento do transporte ferroviário.
\choice do arrefecimento dos ideais absolutistas no Brasil e a disseminação de movimentos nativistas.
\CorrectChoice da emancipação política do Brasil e o seu ingresso na órbita da influência britânica.
\choice da persistência do pacto colonial no Brasil e o seu ingresso no capitalismo monopolista.
\choice do fechamento das fronteiras do Brasil aos estrangeiros e a abertura para as correntes ideológicas revolucionárias europeias.
\end{choices}
